\section{Đánh giá và kết luận}\label{ch:ketluan}

\subsection{Đối chiếu với yêu cầu nộp bài}
\begin{longtable}{@{}L{5.6cm}L{1.6cm}L{6.6cm}@{}}
\caption{Đối chiếu checklist nộp bài của đề}\label{tab:checklist}\\
\toprule
\textbf{Hạng mục} & \textbf{Trạng thái} & \textbf{Vị trí} \\ \midrule \endfirsthead
\toprule \textbf{Hạng mục} & \textbf{Trạng thái} & \textbf{Vị trí} \\ \midrule \endhead
\bottomrule \endfoot
Link repository & Có & \url{https://github.com/tthan-prog/BrewLite-AS2} \\
SRS, đặc tả use case, ma trận truy vết & Có & \code{docs/\allowbreak{}requirements/\allowbreak{}}, \code{docs/\allowbreak{}traceability-matrix.\allowbreak{}md}; mục \ref{ch:yeucau}, \ref{ch:truyvet} \\
Activity, sequence, class, deployment diagram & Có & \code{docs/\allowbreak{}diagrams/\allowbreak{}}, \code{docs/\allowbreak{}architecture/\allowbreak{}}; mục \ref{ch:mohinh}, \ref{ch:kientruc} \\
Mã nguồn frontend, backend, migration & Có & \code{apps/web}, \code{apps/api}, \code{apps/\allowbreak{}api/\allowbreak{}prisma/\allowbreak{}migrations} \\
Unit, integration, end-to-end test & Có & mục \ref{ch:kiemthu} \\
Pipeline CI/CD chạy thành công & Có & \code{.\allowbreak{}github/\allowbreak{}workflows/\allowbreak{}}; CD \code{v1.0.0}, \code{v1.1.0}, \code{v1.1.1} xanh \\
URL production, API, health check & Có & mục \ref{sec:production} \\
Release notes, README, hướng dẫn rollback & Có & \code{docs/release/}, \code{README.md}, \code{docs/\allowbreak{}deploy/\allowbreak{}runbook.\allowbreak{}md} \\
Sprint Review và Retrospective & Có & \code{docs/scrum/}; có phản hồi của Product Owner và góp ý của cả 4 thành viên cho Sprint 1--3 (mục \ref{ch:sprint1}, \ref{ch:sprint2}, \ref{ch:sprint3}) \\
Video demo & Chưa có & Ghi ngoài repository \\
\end{longtable}

\subsection{Kết quả đạt được}
\begin{enumerate}
\item \textbf{Đi đủ năm giai đoạn đúng thứ tự}, có lịch sử thay đổi rõ ràng trên một repository riêng. Mọi thay đổi vào \code{main} đều qua pull request và CI.
\item \textbf{Nghiệp vụ khó được thiết kế từ đầu và có test với database thật:}
  \begin{itemize}
  \item không bán vượt nguyên liệu khi thanh toán đồng thời;
  \item không trừ tiền hai lần;
  \item máy trạng thái đơn;
  \item hoàn kho, hoàn điểm khi hủy.
  \end{itemize}
\item \textbf{Pipeline CD hoàn chỉnh:} từ tag phiên bản đến production, có smoke test kiểm tra đúng phiên bản và rollback bằng tag cũ. Pipeline đã chứng minh giá trị khi phát hiện lỗi ``xanh nhưng vẫn chạy bản cũ'' ở \code{v0.9.1}.
\item \textbf{Sản phẩm chạy được trên production qua HTTPS} với bốn phương thức thanh toán, trong đó VNPay sandbox được xác nhận từ phía VNPay.
\item \textbf{Truy vết đầy đủ} từ FR, NFR đến user story, sơ đồ, mã nguồn và test.
\end{enumerate}

\subsection{Hạn chế}\label{sec:hanche}
\begin{itemize}
\item \textbf{Quy trình review:}
  \begin{itemize}
  \item Ruleset đặt 0 approval, nhiều pull request được merge khi chưa có thành viên khác approve. Mục 4 của Definition of Done chưa được tuân thủ đầy đủ.
  \item Sprint 3 do một người hiện thực.
  \end{itemize}
\item \textbf{Phân bố công việc không đều:} phần lớn mã nguồn Task 5 và Sprint 3 do Scrum Master viết, có công cụ AI hỗ trợ. Các thành viên khác chủ yếu làm sơ đồ, schema và giao diện (bảng \ref{tab:donggop}).
\item \textbf{Thời gian rất ngắn:} mỗi Sprint khoảng một ngày. Daily Scrum và các sự kiện Scrum diễn ra qua tin nhắn, nên khó tạo đủ nhịp ``kiểm tra -- thích nghi''.
\item \textbf{Biên bản bổ sung muộn:} phản hồi của Product Owner và góp ý riêng của các thành viên được điền ngày 17/11/2026, sau khi các Sprint đã kết thúc. Video demo chưa có.
\item \textbf{Kỹ thuật:}
  \begin{itemize}
  \item Test hiệu năng (PERF-T01) mới đo độ trễ tuần tự, chưa có load test nhiều người dùng đồng thời.
  \item JWT không thu hồi được trước khi hết hạn.
  \item Gói Free của Render ngủ khi không có truy cập (lần mở đầu chậm khoảng 50 giây), và database Free có thời hạn sử dụng theo chính sách của Render.
  \item Cổng Ví, Thẻ, QR là mô phỏng; VNPay chỉ dùng sandbox.
  \end{itemize}
\item \textbf{Kiểm thử giao diện:} E2E chưa bắt được các lỗi chỉ xuất hiện ở màn hình rộng hoặc trên image Docker (mục \ref{ch:sprint3}).
\end{itemize}

\subsection{Bài học}
\begin{itemize}
\item \textbf{Tiêu chí chấp nhận cụ thể làm việc viết test dễ hơn nhiều.} Phần lớn tên test case lấy gần như nguyên văn từ tiêu chí chấp nhận trong SRS.
\item \textbf{Nghiệp vụ đồng thời phải được thiết kế ở mức kiến trúc} (ADR 005) và kiểm chứng trên database thật. Mock không phát hiện được lỗi bán vượt kho.
\item \textbf{Smoke test chỉ kiểm tra ``dịch vụ còn sống'' là chưa đủ.} Phải kiểm tra đúng phiên bản vừa triển khai.
\item \textbf{Kiểm thử thủ công trên môi trường thật vẫn cần thiết} dù có nhiều test tự động. Ba lỗi của \code{v1.1.0} chỉ lộ ra trên production.
\item \textbf{Thống nhất hợp đồng API trước khi làm giao diện} giúp ghép frontend và backend ngay lần đầu.
\end{itemize}

\subsection{Hướng phát triển}
\begin{itemize}
\item Bật yêu cầu ít nhất 1 approval trong ruleset; luân phiên người review.
\item Thêm E2E cho đăng xuất nhân viên, manifest có biểu tượng PNG, và chạy E2E trên image Docker.
\item Mở rộng PERF-T01 thành load test đồng thời (ví dụ k6).
\item Giao diện quản trị sản phẩm, khuyến mãi và báo cáo doanh thu (FR-22).
\item Kết nối cổng thanh toán thật sau khi đáp ứng yêu cầu pháp lý (FR-23).
\item Thu hồi phiên đăng nhập (refresh token hoặc danh sách token bị thu hồi).
\end{itemize}

\subsection{Kết luận}
Nhóm đã xây dựng BrewLite theo đúng chuỗi công việc của Bài tập lớn 2: yêu cầu $\rightarrow$ mô hình hóa $\rightarrow$ kiến trúc $\rightarrow$ hai Sprint hiện thực $\rightarrow$ CI/CD và triển khai.

Phiên bản \code{v1.1.1} đang chạy trên production. Luồng nghiệp vụ tối thiểu của đề chạy được và được kiểm chứng bằng 211 test tự động cùng kiểm thử thủ công.

Các hạn chế về quy trình review và phân bố công việc là điểm nhóm cần cải thiện nhất trong các dự án sau.

\bigskip
